%
% ---------- header -----------------------------------------------------------
%
% project       kaneton
%
% license       kaneton
%
% file          /home/mycure/kaneton/view/book/kaneton/kaneton.tex
%
% created       julien quintard   [mon may 14 19:56:45 2007]
% updated       julien quintard   [fri jun  1 15:02:48 2007]
%

%
% path
%

\def\path{../..}

%
% template
%

%
% ---------- header -----------------------------------------------------------
%
% project       kaneton
%
% license       kaneton
%
% file          /home/mycure/kaneton/view/template/book.tex
%
% created       julien quintard   [wed may 16 18:16:50 2007]
% updated       julien quintard   [sun jun  3 12:00:37 2007]
%

%
% class
%

\documentclass[10pt,a4wide]{book}

%
% packages
%

\usepackage[english]{babel}
\usepackage[T1]{fontenc}
\usepackage{a4wide}
\usepackage{fancyheadings}
\usepackage{multicol}
\usepackage{indentfirst}
\usepackage{graphicx}
\usepackage{color}
\usepackage{xcolor}
\usepackage{verbatim}

\usepackage{aeguill}

\usepackage[Lenny]{\path/package/fncychap}

\pagestyle{fancy}

%
% -
%

\renewcommand{\-}{\vspace{\parskip}}

%
% toc
%

\newcommand{\toc}
  {
    \tableofcontents
    \setlength{\footrulewidth}{0.3pt}
    \setlength{\parindent}{0.3cm}
    \setlength{\parskip}{2ex plus 0.5ex minus 0.2ex}
  }

%
% logos
%

\newcommand{\logos}
  {
    \begin{center}
      \includegraphics[scale=0.8]{\path/logo/kaneton.pdf}
    \end{center}
  }

%
% colors
%

\definecolor{functioncolor}{rgb}{0.40,0.00,0.00}
\definecolor{commandcolor}{rgb}{0.00,0.00,0.40}
\definecolor{verbatimcolor}{rgb}{0.00,0.40,0.00}
\definecolor{noticecolor}{rgb}{0.87,0.84,0.02}

%
% function
%

\newcommand\function[3]{
  \begin{tabular}{p{0.2cm}p{13.8cm}}
  & {\color{functioncolor}\textbf{#1}}#2
  \end{tabular}

  \begin{tabular}{p{1cm}p{13cm}}
  & #3
  \end{tabular}}

%
% align
%

\newcommand\align[1]{
  \\ & \hspace{#1}}

%
% argument
%

\newcommand\argument[1]{\textit{#1}}

%
% command
%

\newcommand\command[2]{
  \begin{tabular}{p{0.2cm}p{13.8cm}}
  & {\color{commandcolor}\textbf{#1}}
  \end{tabular}

  \begin{tabular}{p{1cm}p{13cm}}
  & #2
  \end{tabular}}

%
% notice
%

\newcommand\notice[1]{
  {\color{noticecolor}\textbf{Notice}}

  \begin{tabular}{p{0.2cm}p{13.8cm}}
  & #1
  \end{tabular}}

%
% subsubsubsection
%

\newcommand\subsubsubsection[1]{\textbf{#1}}

%
% example
%

\newcommand\example[1]{
  \textit{Example:}

  \begin{tabular}{p{0.2cm}p{13.8cm}}
  & \textit{#1}
  \end{tabular}}

%
% warning XXX
%

\renewcommand{\familydefault}{\sfdefault}

%
% verbatim stuff
%

\makeatletter

\renewcommand{\verbatim@font}
  {\ttfamily\footnotesize\selectfont}

\def\verbatim@processline{
  \hskip15ex{\color{verbatimcolor}\the\verbatim@line}\par
}

\makeatother

%
% header
%

\rhead{}
\rfoot{\scriptsize{The kaneton microkernel project}}

\date{\scriptsize{\today}}


%
% header
%

\lhead{\scriptsize{The kaneton microkernel}}
\rhead{}

%
% title
%

\title{The kaneton microkernel
       \logos}

%
% authors
%

\author{\small{Julien Quintard}}

%
% document
%

\begin{document}

%
% title
%

\maketitle

%
% --------- text --------------------------------------------------------------
%

This document describes the kaneton microkernel research project.

This document should be used by every student willing implement the
kaneton microkernel as well as by people looking for more details on the
kaneton microkernel design and implementation.

All the kaneton documents are available on
the official website
  \footnote{http://www.kaneton.org}.

%
% toc
%

\tableofcontents

%
% chapters
%

%%
%% licence       kaneton licence
%%
%% project       kaneton
%%
%% file          /home/mycure/kaneton/view/papers/kaneton/goals.tex
%%
%% created       matthieu bucchianeri   [mon jan 30 17:09:17 2006]
%% updated       julien quintard   [thu mar  2 01:21:16 2006]
%%

%
% goals
%

\chapter{Goals}

The kaneton microkernel is studied during the EPITA System, Network
and  Security specialization. The project comes with many courses
leading student to learn about architectures and kernels and to design
and develop important parts of a microkernel.

From this fact, kaneton has been designed in a very particular way to
keep the design as simple as possible to be understandable by the
students.

kaneton is not based on any previous kernel code: it is written from
scratch. As we wanted the kernel to be portable on any architecture,
we paid particular attention in dividing the code into two precise
sections: the architecture-independent code and the
architecture-dependent code.

---

We also use terms to design these two different pieces of code, the
\textbf{core} and the \textbf{arch} respectively. These terms are not
quite defined and are not very clear and expressive but we chose them.

---

Designing this portability system, we intended to port the kernel on
many architectures, so we had to keep the architecture-dependent code
as small as possible. Moreover, the only requirement about the
hardware is the minimal amount of memory, which must be at least 32
megabytes.

As we wanted the project to be very pedagogic, kaneton source code is
clear and heavily commented. Each manager has its own directory and
headers. Function names and prototypes are unified and explicit.
Each function is preceded by a header explaining the task it performs
and detailing the different steps.

%
% history
%

\section{History}

%
% year 2006
%

\subsection{Year 2006}

The kaneton project comes with a development environment. Then the
students have to write parts of the microkernel without having to
build a development environment from scratch which takes much time.

Moreover, courses were added to the EPITA specialization curriculum
including microprocessor's architectures, kernels history etc..

%
% year 2005
%

\subsection{Year 2005}

The first complete kaneton project was scheduled from January to December.

%
% year 2004
%

\subsection{Year 2004}

The project was introduced by students willing learn kernel architectures.

%%
%% licence       kaneton licence
%%
%% project       kaneton
%%
%% file          /home/mycure/kaneton/view/papers/kaneton/history.tex
%%
%% created       julien quintard   [sat mar  4 14:15:39 2006]
%% updated       julien quintard   [mon may  8 17:48:32 2006]
%%

%
% history
%

\chapter{History}

In this chapter we will detail the kaneton history from the first
year with low-level programming introduction to the last kaneton
reference implementation.

\newpage

%
% text
%

\textit{Since this chapter is intended to retrace the kaneton history,
  no insurance can be made on emails validity.}

As said in the previous chapter, the kaneton microkernel was introduced
by two students
Julien Quintard
  \footnote{julien.quintard@gmail.com} and
Jean-Pascal Billaud
  \footnote{billau\_j@epita.fr}
during the year 2004.

The kaneton microkernel is currently studied during the
EPITA
  \footnote{http://www.epita.fr}
System, Network and Security specialization curriculum
  \footnote{http://srs.epita.fr}.

During the kaneton history, the project evolved and lectures were added
to the curriculum to make the whole kaneton project more interesting and
understandable by the students.

%
% year 2004
%

\section{Year 2004}

The first year, a low-level programming introduction course was proposed
for the engineering school's first year students.

Then, about fourteen hours courses were made to introduce microprocessor's
external architecture and low-level programming.

The students had to develop small, poor and messy device drivers for the
console and keyboard peripherals. Moreover, a tiny shell was developed
by students so they were able to enter a command to launch a special
kernel action.

The course was a bit chaotic but this first shot was a success.

Therefore, the System, Network and Security specialization students
themselves asked the two students a more evolved project in their
curriculum so they were able to learn about operating systems internals.

%
% year 2005
%

\section{Year 2005}

Then, for the next year, the two students prepared a complete microkernel
design with two complete courses on kernel design and Intel Architecture
programming.

Moreover, other students joined the project:
C\'edric Aubouy
  \footnote{aubouy\_c@epita.fr},
Fabien Le-Mentec
  \footnote{le-men\_f@epitech.net} and
Renaud Lienhart
  \footnote{lienha\_r@epita.fr}.

The project was composed of six steps, from the bootstrap, passing by
the kernel internals including memory management, task management etc..
to the servers with an IDE device driver and finally a FAT file system.

Once again, the whole project was a success but the kaneton people
noticed that the students took much time doing boring work like
filling in header files, dealing with versionning problems, writing
makefiles and shell scripts etc..

Moreover, the courses were too messy and the students had difficulties
to make the relation between kaneton design and microprocessor's
architecture implementation.

Finally, kaneton people started to implement a kaneton microkernel
reference in C language.

%
% year 2006
%

\section{Year 2006}

People joined the project:
Matthieu Bucchianeri
  \footnote{bucchi\_m@epita.fr} and
Renaud Voltz
  \footnote{voltz\_r@epita.fr}.

These students actively contributed to the kaneton implementation.

This year, kaneton people decided to introduce a development environment,
based on the kaneton implementation reference, including everything
necessary to set up a collaborative kernel development.

While, previously, the students had to write the entire microkernel
and servers from scratch, this year, students only had to write precise
parts of the microkernel.

These major modifications in the kaneton project were made to
allow students to explore advanced fields like distributed systems
aspects and/or microkernel's security.

Finally, courses were added to the EPITA System, Network and Security
specialization curriculum including microprocessor's architectures,
kernels history etc..

The kaneton microkernel implementation, in 2006, counted about
\textit{7,000} lines for the core and about \textit{2,000} lines for the
microprocessor's architecture implementation.

Nevertheless, the microkernel is not complete yet, so a full implementation
of about \textit{15,000} lines can be expected which seems correct
considering the kaneton microkernel advanced design.

%%
%% licence       kaneton licence
%%
%% project       kaneton
%%
%% file          /home/mycure/kaneton/view/papers/kaneton/overview.tex
%%
%% created       matthieu bucchianeri   [mon jan 30 17:09:45 2006]
%% updated       julien quintard   [thu mar  2 13:12:22 2006]
%%

%
% overview
%

\chapter{Overview}

XXX ce chapitre va vous aider a reconnaitre les fonctionnalites principale
XXX d'un kernel dans kaneton.

The kaneton microkernel is only the core of an operating system.
Main tasks like hardware drivers or user services are implemented as
\textbf{servers}. So the microkernel only has a few functionalities to
provide:

\begin{itemize}
  \item
    Memory management.
  \item
    Process management.
  \item
    Communication.
  \item
    Events.
\end{itemize}

In this chapter we will describe briefly these tasks and all the
associated managers.

%
% memory management
%

\section{Memory Management}

Handling the memory -- from virtual address space to physical
addressing -- is done by three major managers, the \textbf{as},
\textbf{segment} and \textbf{region} managers.

%
% as
%

\subsection{as}

The address space manager just manages the different address spaces
used by the kaneton tasks.

In kaneton, we call an \textbf{as - address space} a list of memory
locations referenced by a task. Each task has its own address space.

%
% segment
%

\subsection{segment}

The segment manager just manages the segments reserved by
the different kaneton entities including the kernel, the drivers etc..

In kaneton terms a \textbf{segment} is a contiguous area of reserved
physical memory.

%
% region
%

\subsection{region}

The region manager keeps track of regions used to map segments for
each address space reserved on the system.

In kaneton, a \textbf{region} is contiguous area of virtual memory
mapping a segment's part.

%
% process management
%

\section{Process Management}

XXX

%
% communication
%

\section{Communication}

XXX

%
% events
%

\section{Events}

XXX

%
% ---------- header -----------------------------------------------------------
%
% project       kaneton
%
% license       kaneton
%
% file          /Users/francois/k...rnels/future/portability/portability.tex
%
% created       julien quintard   [fri oct 24 17:31:58 2008]
% updated          [mon jan 26 21:58:11 2009]
%

%
% ---------- setup ------------------------------------------------------------
%

%
% path
%

\def\path{../../../..}

%
% template
%

%
% ---------- header -----------------------------------------------------------
%
% project       kaneton
%
% license       kaneton
%
% file          /home/mycure/kaneton/view/template/lecture.tex
%
% created       julien quintard   [wed may 16 18:17:26 2007]
% updated       julien quintard   [sun may 18 23:23:40 2008]
%

%
% class
%

\documentclass[8pt]{beamer}

%
% packages
%

\usepackage{pgf,pgfarrows,pgfnodes,pgfautomata,pgfheaps,pgfshade}
\usepackage[T1]{fontenc}
\usepackage{colortbl}
\usepackage{times}
\usepackage{amsmath,amssymb}
\usepackage{graphics}
\usepackage{graphicx}
\usepackage{color}
\usepackage{xcolor}
\usepackage[english]{babel}
\usepackage{enumerate}
\usepackage[latin1]{inputenc}
\usepackage{verbatim}
\usepackage{aeguill}

%
% style
%

\usepackage{beamerthemesplit}
\setbeamercovered{dynamic}

%
% verbatim stuff
%

\definecolor{verbatimcolor}{rgb}{0.00,0.40,0.00}

\makeatletter

\renewcommand{\verbatim@font}
  {\ttfamily\footnotesize\selectfont}

\def\verbatim@processline{
  {\color{verbatimcolor}\the\verbatim@line}\par
}

\makeatother

%
% -
%

\renewcommand{\-}{\vspace{0.4cm}}

%
% date
%

\date{\today}

%
% logos
%

\pgfdeclareimage[interpolate=true,width=34pt,height=18pt]
                {epita}{\path/logo/epita}
\pgfdeclareimage[interpolate=true,width=49pt,height=18pt]
                {upmc}{\path/logo/upmc}
\pgfdeclareimage[interpolate=true,width=25pt,height=18pt]
                {lse}{\path/logo/lse}

\newcommand{\logos}
  {
    \pgfuseimage{epita}
  }

%
% institute
%

\institute
{
  \inst{1} kaneton microkernel project
}

%
% table of contents at the beginning of each section
%

\AtBeginSection[]
{
  \begin{frame}<beamer>
   \frametitle{Outline}
    \tableofcontents[current]
  \end{frame}
}

%
% table of contents at the beginning of each subsection
%

\AtBeginSubsection[]
{
  \begin{frame}<beamer>
   \frametitle{Outline}
    \tableofcontents[current,currentsubsection]
  \end{frame}
}


%
% title
%

\title{Portability}

%
% document
%

\begin{document}

%
% title frame
%

\begin{frame}
  \titlepage
\end{frame}

%
% outline frame
%

\begin{frame}
  \frametitle{Outline}

  \tableofcontents
\end{frame}

%
% ---------- text -------------------------------------------------------------
%

%
% introduction
%

\section{Introduction}

\begin{frame}
  \frametitle{Introduction}

  Nowadays, a lot of CPU types exist (IA-32, IA-32\_64, IA-64, ARM, ARM, MIPS, PowerPC, SH, m68k, Blackfin, ...) and a lot of machines uses these CPUs, in different ways.

  \-

  The modern operating systems are trying to cope with that, and they tend to support all these platforms, so that a user can use the same system on all the different machines he uses.

  \-

  Portability requires some design in the kernel, to avoid rewriting everything from scratch and maintaining a separate branch for each supported architecture.

  \-

  This course describes what are the differences between machines, and how to design a kernel so it can be ported on several platforms.

\end{frame}

\section{Architecture}
\subsection{Machine architecture vs CPU architecture}

\begin{frame}
  \frametitle{Machine architecture vs CPU architecture}

  From a kernel point of view, there are two kind of architectures on a machine :

  \begin{itemize}
  \item The CPU Architecture
  \item The Machine/Platform Architecture
  \end{itemize}

  \-

  The following slides will explain what makes a CPU architecture specific, what makes a Machine architecture specific, and how a kernel can be designed so it can be as easily as possible ported on several machines.  

\end{frame}

\subsection{Microprocessor architecture}

\begin{frame}
  \frametitle{Microprocessor architecture}

  The architecture of a microprocessor defines :

  \begin{itemize}
  \item Its intructions set
  \item Its registers
  \item Its operational modes/behaviour
  \end{itemize}

  \-

  It's the interface between the software and the CPU itself, so it's basically what will be documented by the manufacturer in the datasheet of the CPU.

\end{frame}

\begin{frame}
  \frametitle{Microprocessor architecture}
  
  Several microprocessors models can share the same architecture.

  \-

  IA32 (commonly called x86) processors are manufactured by Intel, AMD, Via, \ldots

  \-

  ARM designs CPU architectures, manufacturers can use these specifications to make their own CPU.

\end{frame}

\begin{frame}
  \frametitle{Microprocessor architecture extensions}

  Some manufacturers are expanding one CPU architecture to provide more features, but keeping the main CPU architecture as a base.

  \-

  A code that was made to run on the base architecture will work on all the derivatives architectures, but not the opposite.

  \-

  This approach has been used quite a lot, for example, with IA-32 architecture, and the additional instruction sets (MMX, SSE, SSE2, SSE3 from Intel, 3dNow from AMD)

  \-

  A portable kernel would use these features if they are available, but must provide a software alternative in the other case. That way, the kernel can benefit from the performance gain provided by these instructions set, but doesn't depend on their presence to work.

\end{frame}

\begin{frame}
  \frametitle{CPU Architecture classes}

  Some categories have been made to distinguish families of CPU architectures :

  \begin{itemize}
  \item RISC Architecture: Reduced Instruction Set Computer - This family of CPU architectures focus on providing a really basic and simple set of instructions, and they try to optimize each instruction as much as possible, so that even something quite complex, that will require several instructions, could be faster than on a CPU where one single instruction would have been used. Such an architecture contains generally a lot of registers, so that programs can work as much as possible on data stored in registers instead of doing several memory accesses.

  \item CISC Architecture: Complex Instruction Set Computer - This family of CPU architectures focus more on providing a consequent set of instructions, so that some usual operations, that could be done using several elementary operations, are available through a single instruction. For that reason, the instruction set of a CISC CPU contains quite a lot of operations. This kind of CPU actually quite often contains a RISC CPU and a Microcode that describes how to do each instruction.

  \end{itemize}

\end{frame}

\begin{frame}
  \frametitle{Microprocessor internal and external architecture}

  All this was about external architecture.

  \-

  There is another kind of Microprocessor architecture : the internal architecture.

  \-

  This is how the CPU is actually achieving the implementation of the interface required by the external architecture.

  \-

  This is the CPU manufacturers core business, it's what makes the difference between two CPUs that have the same external architecture, like Intel vs AMD, or Via on IA-32 processors, since it's what makes the difference in terms of performance, power consumption, \ldots

\end{frame}

\subsection{Machine architecture}

\begin{frame}
  \frametitle{Machine architecture}
  
  A microprocessor can't work alone. It has interfaces to communicate with other components : memory buses, peripherial buses, interrupts lines, \ldots

  \-

  This provides a way to connect the memory required by the microprocessor, and some peripherials that can be useful in a machine, such as time-sources/clocks, user interfaces, data storage, network interfaces, \ldots

  \-

  The firmware provided in a machine is also something specific to the machine architecture since it provides services to interface with the hardware to the operating system.

  \-

  All these things are external components that are connected to the CPU in a specific way. A machine architecture describes the set of components that are used, and how they are connected together.

\end{frame}

\begin{frame}
  \frametitle{Machine architecture}

  The IBM-PC architecture is the most known architecture that uses an IA-32 CPU, but there are actually some other. SGI, for instance, made a machine based on an IA-32 CPU, but that was different from the IBM-PC architecture (SGI 320) since it didn't have a BIOS, but a firmware based on their own ARCS firmware they used on their other MIPS machines.

  \-

  More recently, Apple released their own architecture based on IA-32, where they also don't use the BIOS, but instead provide a firmware called EFI.

  \-

  Some other CPU architectures are used in several different machine architectures. For example, the MIPS R5000 CPU was used in the SGI Indy architecture, and in the Sony Playstation 2 architecture. These two machines use the same CPU architecture but they are totally different. An OS written for the SGI Indy won't work on the Sony Playstation 2.

  \-

  Almost every single Windows mobile based smartphone is a different Machine architecture, although they all share the ARM CPU architecture.

\end{frame}

\section{Kernel splitting}
\subsection{Independant part}

\begin{frame}
  \frametitle{Independant part}

  The independant part of a kernel consists in all the code that will work on every machine.

  \-

  It basically consists of all the code handling the kernel's concepts. For example, a task is a high level concept that we can find in almost all kernels. A portable kernel will be able to manage tasks on all the architectures it supports. Some of the task management done by the kernel consists in the manipulation of structures describing the task in a generic way. All this code does not rely on machines specific features, the same code will run and work on all machines. This is an independant code.

  \-

  Of course, creating a task will certainly require some specific actions as well. The goal of a portable kernel will be to isolate as much as possible those parts of the code, and to put as much things as possible in the independant part, to avoid redundancy in the dependant part.

  \-

  No ASM code can be found in the independant part since ASM is by definition dependant of the CPU architecture.

\end{frame}

\subsection{Dependant part}

\begin{frame}
  \frametitle{Dependant part}

  The dependant part of a kernel contains all the code that is aimed to support a specific feature of a CPU, a machine, or a specific hardware.

  \-

  Writing entries in the MMU cache is something specific to the CPU. It's generally done using some assembly instructions that are different from one CPU architecture to another. That's why it's a code that depends on the CPU architecture.

  \-

  Setting up the platform to trigger an interrupt on a regular basis is something specific to the Machine architecture. In general, a machine contains a dedicated hardware, connected to a CPU interrupt line. That's why configuring it is done by some code that depends on the Machine architecture.

  \-

  These examples are not always true : if a CPU embeds its own clock source that can trigger an interrupt internally, then the code to configure it is not Machine architecture dependant, but it is CPU architecture dependant.

\end{frame}

\subsection{Board support package (BSP)}

\begin{frame}
  \frametitle{Dependant part}

  Some kernels, especially the kernels for embedded operating systems, introduced the Board Support Package (BSP) concept.
  
  \-

  For example, Windows CE, Microsoft's kernel for embedded systems, is working on 3 CPU architectures : IA-32, ARM9, SuperH, but it runs on a lot of devices with a lot of various hardware devices. Each device has its own machine architecture, and its own set of peripherials. For that reason, each platform requires some specific code to make the kernel to work on it.

\end{frame}


\section{Kaneton portability}

\begin{frame}
  \frametitle{Kaneton Portability}

  The Kaneton microkernel is designed to be portable. For that reason, the code is splitted in several sections :

  \begin{itemize}
  \item Core
  \item Machine
  \begin{itemize}
  \item Architecture
  \item Platform
  \item Glue
  \end{itemize}
  \end{itemize}

\end{frame}

\subsection{Core}

\begin{frame}
  \frametitle{Kaneton Core}

  In the Kaneton design, the machine independant code is called Core.

  \-

  This section of code contains all the kernel code which is not specific to a machine. It contains the high level functions and code that are used for all the kernel internal operations, such as :
  \begin{itemize}
  \item Allocate physical memory
  \item Allocate and map virtual memory
  \item Create a new task
  \item Send a message between two tasks
  \item \ldots
  \end{itemize}

  Most of these operations don't require to do any CPU specific operations, but some of them do. For example, mapping virtual memory to physical memory requires to configure the MMU through the caches, this is something different for each CPU. For that reason, every function in the Core module of Kaneton will make a call to a potential machine dependant code.

  \-
  
  This is achieved through a macro called machine\_call.

\end{frame}

\subsection{Machine}

\begin{frame}
  \frametitle{Kaneton Architecture}
  
  This section of the Kaneton code contains all the code specific to the CPU architecture.

  \-

  This is mainly the code that will require to use specific assembly calls (inline assembly) or the code that works on data structures imposed by the CPU.

\end{frame}

\begin{frame}
  \frametitle{Kaneton Platform}
  
  This section of the Kaneton code contains all the code specific to the machine architecture.

  \-

  This is the code that will handle some peripherials attached to the CPU in a specific architecture, and that are used by the kernel itself.

\end{frame}

\begin{frame}
  \frametitle{Kaneton Glue}

  In the Kaneton design, the independant code (Core) calls the dependant code.

  \-

  To make it possible, the dependant code must have a generic interface, so that the Core code doesn't contain specifically any call to a function specific to an architecture.

  \-

  In Kaneton, the independant code and the dependant code functions share exactly the same prototypes, so that the Core code can call the Machine code in a generic way.

  \-

  For that reason, a wrapper was required, to implement the correct interface, that would then call the Architecture and the Machine code accordingly. This is the role of Glue.

\end{frame}

\section{Conclusion}

\begin{frame}
  \frametitle{Conclusion}

  In order to make a portable kernel, one must well distinguish what code can be reused for every machine, and what code is specific to the machine he is working on.

  \-

  This problem has now been addressed in most of the modern operating systems kernels (Linux, BSD, Windows CE, \ldots) with more or less style, but it's generally much better when a kernel has been designed to be portable from the beginning.

  \-

  This portability is very important nowadays, since the hardware is evolving fast. Making a new kernel without thinking about portability is a bad strategy since it makes it dependant on the durability of the architecture it is being made for.

  
\end{frame}


%
% bibliography
%

\begin{frame}[allowframebreaks]
  \frametitle{Bibliography}

  \bibliographystyle{amsplain}
  \bibliography{\path/bibliography/bibliography}
\end{frame}

\end{document}

% core
%%
%% licence       kaneton licence
%%
%% project       kaneton
%%
%% file          /home/mycure/kaneton/view/papers/kaneton/architectures.tex
%%
%% created       julien quintard   [sun apr 23 17:08:41 2006]
%% updated       julien quintard   [sun apr 23 17:08:41 2006]
%%

%
% architectures
%

\chapter{Architectures}

This chapter will overview the different architectures supported by
the kaneton microkernel reference.

\newpage

%
% text
%

The kaneton microkernel references currently supports the following
architectures.

Notice that, sometimes, the kaneton microkernel can support multiples
\textit{architectures} of a single microprocessor's architecture.
These multiples \textit{architectures} generally exploit different
facilities of the same microprocessor's architecture.

There is a paper specific to every supported architecture. These papers
are available on the kaneton official website
  \footnote{http://www.kaneton.org}.

%
% intel architecture 32-bit
%

\section{Intel Architecture 32-bit}

The kaneton microkernel reference was first developed on the Intel Architecture
32-bit since this architecture is very popular and cheap.

The following architecture implementations are supported by the kaneton
microkernel.

%
% ia32-virtual
%

\subsection{ia32-virtual}

This architecture uses a flat segmentation model with paging enabled
allowing multiple virtual address space and so virtual address space
protections.

Nevertheless, this architecture implementation is kept as simple as possible
since it is used as the basis of the operating system courses.

One of the particularities of this architecture is that the core has its
own address space. Then, each time a system call occurs, an address space
switch is performed to work in the core's address space.

This particularity leads to bad performances since each time an address
space switch occurs, the microprocessor's caches are flushed.

%
% ia32-optimized
%

\subsection{ia32-optimized}

This architecture also relies on the Intel Architecture 32-bit but uses
every optimisations provided by the microprocessor's architecture.

\notice{This architecture is not implemented yet.}

%
% ---------- header -----------------------------------------------------------
%
% project       kaneton
%
% license       kaneton
%
% file          /home/mycure/kaneton/view/book/core/licenses.tex
%
% created       julien quintard   [tue may 12 18:35:31 2009]
% updated       julien quintard   [tue may 12 18:35:33 2009]
%

%
% ---------- licenses ---------------------------------------------------------
%

\chapter{Licenses}
\label{chapter:licenses}

In this chapter we will details the kaneton-related licenses.

\newpage

%
% ---------- text -------------------------------------------------------------
%

The kaneton project might be considered as an open-project since source
code is provided.

Nevertheless this is not the case as this project is used as material for
operating system courses.

Therefore, people implementing the kaneton microkernel should not
make their source code available. To avoid problems, especially students
cheating, kaneton people decided to use a kaneton-specific license
forbidding source code distribution.

The \name{kaneton license} is based on a more generic license, the
\name{pedagogical licence}.

The next sections will contain these licenses' descriptions.

%
% pedagogical license
%

\section{Pedagogical License}

\textbf{TERMS AND CONDITIONS FOR COPYING, DISTRIBUTION AND MODIFICATION}

\begin{enumerate}
  \item
    This License applies to any program or other work which contains a
    notice placed by the copyright holder saying it may be distributed
    under the terms of this License. The "Program", below, refers
    to any such program or work, and a "work based on the Program"
    means either the Program or any derivative work under copyright law:
    that is to say, a work containing the Program or a portion of it,
    either verbatim or with modifications and/or translated into another
    language. (Hereinafter, translation is included without limitation in
    the term "modification".) Each licensee is addressed as "you".
  \item
    You must not copy or distribute copies of the Program's source
    code, object code or executable form without explicit authorization
    from the maintainers.

    If you have this authorization, you must conspicuously and
    appropriately publish on each copy an appropriate copyright notice
    and disclaimer of warranty; keep intact all the notices that refer
    to this License and to the absence of any warranty; and give any
    other recipients of the Program a copy of this License along with
    the Program.
  \item
    You may modify your copy or copies of the Program or any portion of
    it, thus forming a work based on the Program, provided that you
    meet all of these conditions :

    \begin{enumerate}
      \item
	You must not publish your work without explicit authorization
	from the maintainers.
      \item
	You must send to the maintainers any work that in whole or in
	part contains or is derived from the Program or any part thereof.
      \item
	You must cause any work that you send, that in whole or in
	part contains or is derived from the Program or any part thereof,
	to be licensed as a whole under the terms of this License.
      \item
	You must cause the modified files to carry prominent notices
	stating that you changed the files and the date of any change.
      \item
	If the modified program normally reads commands interactively
	when run, you must cause it, when started running for such
	interactive use in the most ordinary way, to print or display an
	announcement including an appropriate copyright notice and a
	notice that there is no warranty (or else, saying that you
	provide a warranty) and telling the user how to view a copy
	of this License. (Exception: if the Program itself is interactive
	but does not normally print such an announcement, your work based
	on the Program is not required to print an announcement.)
    \end{enumerate}
  \item
    Access to the Program's source is granted if either:

    \begin{enumerate}
      \item
	You want to make the Program evolve
      \item
	You have pedagogical goals
    \end{enumerate}

    These requirements apply to the modified work as a whole. If
    identifiable sections of that work are not derived from the Program,
    and can be reasonably considered independent and separate works in
    themselves, then this License, and its terms, do not apply to those
    sections when you distribute them as separate works. But when you
    distribute the same sections as part of a whole which is a work based
    on the Program, the distribution of the whole must be on the terms of
    this License, whose permissions for other licensees extend to the
    entire whole, and thus to each and every part regardless of who wrote
    it.

    Thus, it is not the intent of this section to claim rights or contest
    your rights to work written entirely by you; rather, the intent is to
    exercise the right to control the distribution of derivative or
    collective works based on the Program.

    In addition, mere aggregation of another work not based on the Program
    with the Program (or with a work based on the Program) on a volume of
    a storage or distribution medium does not bring the other work under
    the scope of this License.
  \item
    You may not copy, modify, sublicense, or distribute the Program
    except as expressly provided under this License. Any attempt
    otherwise to copy, modify, sublicense or distribute the Program is
    void, and will automatically terminate your rights under this
    License. However, parties who have received copies, or rights,
    from you under this License will not have their licenses
    terminated so long as such parties remain in full compliance.
  \item
    You are not required to accept this License, since you have not
    signed it. However, nothing else grants you permission to modify
    or distribute the Program or its derivative works. These actions
    are prohibited by law if you do not accept this
    License. Therefore, by modifying or distributing the Program (or
    any work based on the Program), you indicate your acceptance of
    this License to do so, and all its terms and conditions for
    copying, distributing or modifying the Program or works based on
    it.
  \item
    Each time you redistribute the Program (or any work based on the
    Program), the recipient automatically receives a license from the
    original licensor to copy, distribute or modify the Program subject
    to these terms and conditions. You may not impose any further
    restrictions on the recipients' exercise of the rights granted
    herein. You are not responsible for enforcing compliance by third
    parties to this License.
  \item
    If the distribution and/or use of the Program is restricted in
    certain countries either by patents or by copyrighted interfaces,
    the original copyright holder who places the Program under this
    License may add an explicit geographical distribution limitation
    excluding those countries, so that distribution is permitted only
    in or among countries not thus excluded. In such case, this License
    incorporates the limitation as if written in the body of this
    License.
  \item
    We may publish revised and/or new versions of this License from
    time to time. It may evolve considering new contributors
    needs. Contact us if you have any request. Such new versions will
    be similar in spirit to the present version, but may differ in
    detail to address new problems or concerns.
  \item
    If you wish to incorporate parts of the Program into other free
    programs whose distribution conditions are different, write to the
    author to ask for permission, we sometimes make exceptions for this.
  \item
    THERE IS NO WARRANTY FOR THE PROGRAM, TO THE EXTENT PERMITTED BY
    APPLICABLE LAW. EXCEPT WHEN OTHERWISE STATED IN WRITING THE
    COPYRIGHT HOLDERS AND/OR OTHER PARTIES PROVIDE THE PROGRAM "AS IS"
    WITHOUT WARRANTY OF ANY KIND, EITHER EXPRESSED OR IMPLIED,
    INCLUDING, BUT NOT LIMITED TO, THE IMPLIED WARRANTIES OF
    MERCHANTABILITY AND FITNESS FOR A PARTICULAR PURPOSE. THE ENTIRE
    RISK AS TO THE QUALITY AND PERFORMANCE OF THE PROGRAM IS WITH
    YOU. SHOULD THE PROGRAM PROVE DEFECTIVE, YOU ASSUME THE COST OF
    ALL NECESSARY SERVICING, REPAIR OR CORRECTION.
  \item
    IN NO EVENT UNLESS REQUIRED BY APPLICABLE LAW OR AGREED TO IN
    WRITING WILL ANY COPYRIGHT HOLDER, OR ANY OTHER PARTY WHO MAY
    MODIFY AND/OR REDISTRIBUTE THE PROGRAM AS PERMITTED ABOVE, BE
    LIABLE TO YOU FOR DAMAGES, INCLUDING ANY GENERAL, SPECIAL,
    INCIDENTAL OR CONSEQUENTIAL DAMAGES ARISING OUT OF THE USE OR
    INABILITY TO USE THE PROGRAM (INCLUDING BUT NOT LIMITED TO LOSS OF
    DATA OR DATA BEING RENDERED INACCURATE OR LOSSES SUSTAINED BY YOU
    OR THIRD PARTIES OR A FAILURE OF THE PROGRAM TO OPERATE WITH ANY
    OTHER PROGRAMS), EVEN IF SUCH HOLDER OR OTHER PARTY HAS BEEN
    ADVISED OF THE POSSIBILITY OF SUCH DAMAGES.
\end{enumerate}

%
% kaneton license
%

\section{kaneton License}

\textbf{TERMS AND CONDITIONS FOR COPYING, DISTRIBUTION AND MODIFICATION}

\begin{enumerate}
  \item
    This licence is nothing more than a link to the pedagogical licence.
  \item
    Any program under the kaneton licence is in fact under the terms and
    conditions of the pedagogical licence.
\end{enumerate}

%%
%% licence       kaneton licence
%%
%% project       kaneton
%%
%% file          /home/mycure/kaneton/view/papers/kaneton/bibliography.tex
%%
%% created       julien quintard   [mon may  8 18:35:35 2006]
%% updated       julien quintard   [mon may  8 20:38:56 2006]
%%

%
% bibliograpy
%

\chapter{Bibliography}

This chapter contains the bibliography.

%
% text
%

\begin{thebibliography}{0}
  \bibitem{AST-SCO}
    \textbf{Structured Computer Organization};
    by
    \textit{Andrew S. Tanenbaum}
  \bibitem{AST-CN}
    \textbf{Computer Networks};
    by
    \textit{Andrew S. Tanenbaum}
  \bibitem{AST-OSDI}
    \textbf{Operating Systems: Design and Implementation};
    by
    \textit{Andrew S. Tanenbaum, Albert S Woodhull}
  \bibitem{AST-MOS}
    \textbf{Modern Operating Systems};
    by
    \textit{Andrew S. Tanenbaum}
  \bibitem{AST-DOS}
    \textbf{Distributed Operating Systems};
    by
    \textit{Andrew S. Tanenbaum}
  \bibitem{AST-DSPP}
    \textbf{Distributed Systems: Principles and Paradigms};
    by
    \textit{Andrew S. Tanenbaum, Maarten van Steen}
  \bibitem{NAL-DA}
    \textbf{Distributed Algorithms};
    by
    \textit{Nancy A. Lynch}
\end{thebibliography}


\end{document}

--

overview: kernel = arch + generic :: data-struct + proc

--

bouger les figures dans view/figures
XXX dans book::kaneton, archi supportees, core, glue, machdep system, syscall
    interface, dans book::core, un fichier par manager.
